
En este apartado abordaremos la parte correspondiente a la práctica 2 de la asignatura.
Se puede observar en la siguiente figura el Diagrama de Flujo de Datos (DFD) 
de nivel 1 del subsistema de tendencias. En este diagrama se representan los diferentes 
procesos, almacenes de datos, agentes externos y flujos de datos que componen el subsistema.

\begin{figure}[h!]
    \centering
    \shorthandoff{<>}
    \resizebox{\textwidth}{!}{%
    \begin{tikzpicture}[node distance=1.8cm and 4cm, auto]

        % === Agentes ===
        \node[agente, minimum width=4cm] (Usuario) {Usuario};
        \node[agente, below=of Usuario, yshift=-11cm] (Admin) {Admin};

        % === Almacén de datos ===
        \node[almacen, right=of Usuario, xshift=11.5cm, minimum width=5.5cm] (Hashtag) {Hashtag};

        % === Procesos ===
        \node[proceso, right=of Usuario, xshift=3cm] (RF21) {RF2.1: Crear o mencionar hashtag};
        \node[proceso, below=of RF21, yshift=-0.5cm] (RF22) {RF2.2: Listar tendencias};
        \node[proceso, below=of RF22, yshift=-0.5cm] (RF24) {RF2.4: Mostrar categoría ordenada};
        \node[proceso, below=of RF24, yshift=-0.5cm] (RF23) {RF2.3: Asignar categoría a un hashtag};
        \node[proceso, below=of RF23, yshift=-0.5cm] (RF25) {RF2.5: Eliminar tendencia};

        % === Flujo de datos ===

        % Usuario → RF2.1
        \draw[flujo, ->] (Usuario.east) -- node[above] {RDE 2.1} (RF21.west);

        % RF2.1 → Hashtag
        \draw[flujo, ->] (RF21.east) -- node[above] {RDW 2.1} (Hashtag.west);

        % Usuario → RF2.2
        \draw[flujo, ->] (Usuario.south) ++(1.2,0) |- node[pos=0.8, above left] {Consulta tendencias} (RF22.west);

        % RF2.2 → Usuario
        \draw[flujo, ->] (RF22.north) -- node[left] {RDS 2.2} (Usuario.south east);

        % Hashtag → RF2.2
        \draw[flujo, ->] (Hashtag.west) -- node[left] {RDR 2.2} (RF22.east);

        % Admin → RF2.3
        \draw[flujo, ->] (Admin.east) -- node[left] {RDE 2.3} (RF23.west);

        % RF2.3 → Hashtag
        \draw[flujo, ->] (RF23.east) -- node[right] {RDW 2.3} (Hashtag.south);

        % Hashtag → RF2.3
        \draw[flujo, ->] (Hashtag.south) ++(1,0) |- node[above left] {RDR 2.3} (RF23.south east);

        % Usuario → RF2.4
        \draw[flujo, ->] (Usuario.south) ++(0,-1.2) |- node[pos=0.8, above left] {RDE 2.4} (RF24.west);

        % Hashtag → RF2.4
        \draw[flujo, ->] (Hashtag.south) -- node[left] {RDR 2.4} (RF24.east);

        % Admin → RF2.5
        \draw[flujo, ->] (Admin.east) -- node[above] {RDE 2.5} (RF25.west);

        % RF2.5 → Hashtag
        \draw[flujo, ->] (RF25.east) |- ++(1.5,0) -| node[pos=0.8, above left] {RDW 2.5} (Hashtag.south east);


        % RF2.4 → Usuario
        \draw[flujo, ->] (RF24.west) |- ++(-2,-1) -| node[pos=0.6, left] {RDS 2.4} (Usuario.south west);

    \end{tikzpicture}
    }
    \caption{DFD1 — Subsistema de tendencias}
    \label{fig:dfd_tendencias}
\end{figure}
